\section*{अध्याय \ref{ch:b2:weather-systems} --- मौसम प्रणालियाँ}

\begin{solution}{exo:b2:weather-systems:1}
(क) $f = 2 \times 7.292 \times 10^{-5} \times \sin\ang{50} = \qty{1.117e-4}{s^{-1}}$। (ख)
$\bar T = \qty{278}{K}$ के साथ $\tan\alpha = 1.117 \times 10^{-4} \times 278 \times 15/(9.81 \times 6)
= 0.00791$, लगभग 126 में 1। (ग) $1/0.00791 = 126$: वाताग्र से लगभग \qty{126}{km} पीछे। (घ) ढलान
तापमान-अंतर के व्युत्क्रमानुपाती है: वह आधी हो जाती है, 253 में 1।
\end{solution}

\begin{solution}{exo:b2:weather-systems:2}
(क) $6 \times 150 = \qty{900}{km}$ आगे। (ख) $900/25 = \qty{36}{h}$, डेढ़ दिन। (ग) पहले
ऊँचे पक्षाभ, फिर पक्षाभ स्तरी, फिर मध्य स्तरी; मेघ घना और नीचा होता जाता है, और वाताग्र के निकट आने पर
निचले, घने मेघ से वर्षा शुरू हो जाती है।
\end{solution}

\begin{solution}{exo:b2:weather-systems:3}
प्रतीक उस दिशा में इंगित करते हैं जिधर वाताग्र चलता है: शीत वाताग्र के आगे की अधिक गर्म वायु में, या उष्ण वाताग्र के आगे की
अधिक ठंडी वायु में। (क) शीत वाताग्र: ठंडी वायु उसके उत्तर-पश्चिम में, उसके पीछे, है; वह दक्षिण-पूर्व की ओर
चलता है। (ख) उष्ण वाताग्र: ठंडी वायु उसके उत्तर में, उसके आगे, है; वह उत्तर की ओर चलता है।
(ग) \omterm{def:b2:weather-systems:fronts}{अधिविष्ट वाताग्र}, जो निम्न दाब को उस त्रिक बिंदु से जोड़ता है जहाँ से शीत और उष्ण वाताग्र
शुरू होते हैं; उसके दोनों ओर ठंडी वायु है, गर्म वायु ऊपर उठा दी गई है; वह उस ओर चलता है जिधर उसके
प्रतीक इंगित करते हैं। (घ) \omterm{def:b2:weather-systems:fronts}{स्थिर वाताग्र}: त्रिभुज दक्षिण की ओर गर्म वायु में इंगित करते हैं, इसलिए ठंडी
वायु उत्तर में है; वह मुश्किल से चलता है।
\end{solution}

\begin{solution}{exo:b2:weather-systems:4}
(क) $\beta = 2 \times 7.292 \times 10^{-5} \times \cos\ang{45}/(6.371 \times 10^6) =
\qty{1.62e-11}{m^{-1}.s^{-1}}$। (ख) \qty{4000}{km} के लिए $k = \qty{1.571e-6}{m^{-1}}$, $\beta/k^2 =
\qty{6.6}{m/s}$ और $c = 12 - 6.6 = \qty{5.4}{m/s}$; \qty{8000}{km} के लिए $\beta/k^2 =
\qty{26.3}{m/s}$ और $c = \qty{-14.3}{m/s}$। (ग) छोटी तरंग पूर्व की ओर चलती है, \qty{470}{km} प्रति दिन;
लंबी तरंग पश्चिम की ओर चलती है, \qty{1240}{km} प्रति दिन, पवन के विरुद्ध।
\end{solution}

\begin{solution}{exo:b2:weather-systems:5}
(क) $c = 0$ जब $U = \beta/k^2$, अतः $k = \sqrt{\beta/U}$ और $\lambda = 2\pi/k =
2\pi\sqrt{U/\beta}$। (ख) $2\pi\sqrt{12/(1.62 \times 10^{-11})} = \qty{5410}{km}$। (ग) \ang{45} का
अक्षांश-वृत्त $2\pi \times 6371 \times \cos\ang{45} = \qty{28300}{km}$ का है: लगभग 5.2
तरंगें, इसलिए व्यवहार में पाँच। (घ) $\lambda$ $\sqrt U$ के अनुपात में बढ़ता है: अधिक प्रबल पछुआ पवन
स्थिर तरंगों को लंबा करती है, और उनमें से कम गोलार्ध के चारों ओर समाती हैं।
\end{solution}

\begin{solution}{exo:b2:weather-systems:6}
(क) $C_+ = (\sqrt{72}, \sqrt{72}, 27)$ पर: $\dot x = 10(\sqrt{72} - \sqrt{72}) = 0$, $\dot y =
\sqrt{72}(28 - 27) - \sqrt{72} = 0$, $\dot z = 72 - \tfrac83 \times 27 = 72 - 72 = 0$।
(ख) $\mathrm{e}^{-13.67} = 1.16 \times 10^{-6}$: आयतन एक समय-इकाई में लगभग दस लाख गुना
सिकुड़ता है। (ग) कोई आयतन तब सिकुड़ सकता है जब वह एक दिशा में खिंचता है और बाकी दिशाओं में कहीं अधिक
प्रबलता से दबता है: खिंची हुई दिशा के अनुदिश दूरी बढ़ती है, और दो निकटवर्ती
प्रक्षेप-पथ उसके अनुदिश अलग होते हैं, जबकि उनके बीच का आयतन लुप्त हो जाता है।
\end{solution}

\begin{solution}{exo:b2:weather-systems:7}
(क) $\ln 2/0.906 = 0.765$ समय-इकाई। (ख) $\ln(10/10^{-8})/0.906 = \ln(10^9)/0.906 = 22.9$
समय-इकाई। (ग) $\ln 1000/0.906 = 7.6$ समय-इकाई: आरंभिक अवस्था में हज़ार गुना सुधार
(ख) में पाए गए समय का केवल एक-तिहाई देता है, क्योंकि $\ln 10^3/\ln 10^9 = 1/3$।
\end{solution}

\begin{solution}{exo:b2:weather-systems:8}
(क) $(303 - 195)/195 = 0.554$ और $V_{\max} = \sqrt{0.9 \times 0.554 \times 15\,000} =
\qty{86.5}{m/s}$। (ख) श्रेणी 5 के \qty{70}{m/s} से काफ़ी ऊपर, उस \qty{86}{m/s} के निकट जिसे
हाल के पाँच तूफ़ानों ने पार किया है। (ग) तूफ़ान शायद ही कभी देर तक उस स्थायी अवस्था में रहते हैं जो सिद्धांत
मानता है; ऊर्ध्वाधर पवन-अपरूपण, तूफ़ान में खींची गई शुष्क वायु और उसके चारों ओर उतरती वायु, जिन्हें
सिद्धांत छोड़ देता है, उन्हें दुर्बल करते हैं।
\end{solution}

\begin{solution}{exo:b2:weather-systems:9}
(क) $(303 - 195)/303 = 0.356$। (ख) $(303 - 195)/195 = 0.554$; अनुपात $0.554/0.356 =
303/195 = 1.554$ है। (ग) कार्नो इंजन में कार्य इंजन से बाहर चला जाता है। उष्णकटिबंधीय चक्रवात में
कार्य समुद्र-सतह के निकट घर्षण से क्षयित होता है, और वह ऊष्मा भीतर आती वायु को
तापमान $T_s$ पर गर्म करती है, जहाँ इंजन अपनी ऊष्मा लेता है: उसका फिर से उपयोग होता है, और इंजन ऐसे व्यवहार करता है
मानो वह गुणक $T_s/T_o$ से अधिक दक्ष हो।
\end{solution}

\begin{solution}{exo:b2:weather-systems:10}
(क) \cref{prop:b2:weather-systems:slope} की बिना सन्निकटन वाली व्युत्पत्ति देती है
$\mathrm{d}z/\mathrm{d}x = f(\rho_c v_c - \rho_w v_w)/[g(\rho_c - \rho_w)]$। हल्के तरल पर
टिका हुआ अधिक घना तरल अस्थिर होता है: ज़रा-सा विक्षोभ उसे पलट देता है। विरामावस्था के किसी वाताग्र में
इसलिए ठंडी, अधिक घनी वायु को नीचे ही होना चाहिए, जो परिभाषा की फन्नी है; तब
$\rho_c > \rho_w$, पृष्ठ ठंडी वायु की ओर उठता है, और $\rho_c v_c > \rho_w v_w$।
(ख) $\zeta = \partial v/\partial x = 10/(50 \times 10^3) = \qty{2e-4}{s^{-1}}$, लगभग $1.9f$, जहाँ
$f = \qty{1.03e-4}{s^{-1}}$। (ग) उत्तरी गोलार्ध में धनात्मक: चक्रवाती, और स्वयं कोरिऑलिस प्राचल जितनी
बड़ी; वाताग्र प्रबल चक्रवाती भ्रमिलता की पट्टियाँ हैं।
\end{solution}

\begin{solution}{exo:b2:weather-systems:11}
(क) $\log_2(6/0.5) = \log_2 12 = 3.58$ दोगुनीकरण, $3.58 \times 2.4 = 8.6$~दिन। (ख) पूर्वानुमान के
एक और दिन के लिए त्रुटि के बेकार होने से पहले एक दोगुना-काल कम वृद्धि चाहिए:
आरंभिक त्रुटि को हर दशक $2^{1/2.4} = 1.33$ गुना घटना होगा। (ग) दस गुना कमी
$\log_2 10 = 3.32$ दोगुनीकरण है; यदि उससे केवल 4 से 5 दिन मिलते हैं, तो छोटी त्रुटियाँ $4/3.32 = 1.2$
से $5/3.32 = 1.5$ दिनों में दोगुनी होती हैं: छोटी त्रुटियाँ 1980 के दशक के पूर्वानुमानों में मापी गई बड़ी त्रुटियों से
तेज़ बढ़ती हैं, छोटे पैमानों पर संवहन में मुक्त होने वाली गुप्त ऊष्मा के माध्यम से।
\end{solution}

\begin{solution}{exo:b2:weather-systems:12}
(क) \qty{30}{\celsius} पर $e_s = \qty{43.5}{hPa}$ और $q_s^* = \qty{27.1}{g/kg}$; \qty{29}{\celsius} पर
वायु के लिए $e = 0.8 \times 41.0 = \qty{32.8}{hPa}$ और $q_a = \qty{20.4}{g/kg}$;
$h_o^* - h^* = 1004.5 \times 1 + 2.501 \times 10^6 \times 0.0067 = \qty{17.8}{kJ/kg}$, जबकि पहले
\qty{16.8}{kJ/kg}। (ख) $(303.15 - 200)/200 = 0.516$ और $V_{\max} = \sqrt{0.9 \times 0.516
\times 17\,800} = \qty{90.9}{m/s}$, जबकि पहले \qty{87.9}{m/s}। (ग) प्रति केल्विन लगभग \qty{3.4}{\percent}।
असाम्य \qty{6}{\percent} बढ़ता है, जो गति में \qty{3}{\percent} जोड़ता है;
गुणक $(T_s - T_o)/T_o$ \qty{1}{\percent} बढ़ता है, जो \qty{0.5}{\percent} जोड़ता है:
ईंधन, यानी वह नमी जो अधिक गर्म समुद्र दे सकता है, $e_s$ की चरघातांकी वृद्धि के माध्यम से
सबसे अधिक योगदान देता है।
\end{solution}

\begin{solution}{pb:b2:weather-systems:1}
\textbf{1.} $(480 - 390)/6 = \qty{15}{km/h}$।
\textbf{2.} 12 UTC के $390/15 = \qty{26}{h}$ बाद: अगले दिन 14 UTC पर।
\textbf{3.} $f = \qty{1.117e-4}{s^{-1}}$।
\textbf{4.} $\bar T = \qty{281}{K}$ के साथ $\tan\alpha = 1.117 \times 10^{-4} \times 281 \times
16/(9.81 \times 8) = 0.0064$, लगभग 156 में 1।
\textbf{5.} तब वाताग्र कस्बे से $6 \times 15 = \qty{90}{km}$ आगे निकल चुका होता है; पृष्ठ उसके ऊपर
$90\,000 \times 0.0064 = \qty{576}{m}$ पर है।
\textbf{6.} वाताग्र से $1500/0.0064 = \qty{234}{km}$ पीछे, जो गुज़रने के $234/15 = \qty{15.6}{h}$
बाद पहुँचता है।
\textbf{7.} कपासी और कपासी वर्षी मेघ, बौछारें, शायद मेघगर्जन, झोंके और तापमान में अचानक
गिरावट; दाब, जो वाताग्र के निकट आते समय गिरता रहता है, उसके गुज़रते ही बढ़ने लगता है; फिर
ठंडी, अधिक साफ़ वायु।
\textbf{8.} $\beta = 2 \times 7.292 \times 10^{-5} \times \cos\ang{50}/(6.371 \times 10^6) =
\qty{1.47e-11}{m^{-1}.s^{-1}}$।
\textbf{9.} $k = 2\pi/(5 \times 10^6) = \qty{1.257e-6}{m^{-1}}$ और $\beta/k^2 = \qty{9.3}{m/s}$।
\textbf{10.} $c = 20 - 9.3 = \qty{10.7}{m/s}$, पूर्व की ओर।
\textbf{11.} $10.7 \times 86\,400 = \qty{924}{km}$ पूर्व की ओर।
\textbf{12.} $\lambda_s = 2\pi\sqrt{20/(1.47 \times 10^{-11})} = \qty{7330}{km}$; \ang{50} का
अक्षांश-वृत्त $2\pi \times 6371 \times \cos\ang{50} = \qty{25700}{km}$ का है, ऐसी लगभग 3.5
तरंगें।
\textbf{13.} उसके आगे, उसकी पूर्वी ओर, जहाँ छोटी तरंगें उठती वायु उत्पन्न करती हैं और
सतही दाब गिरता है।
\textbf{14.} $\lambda = \ln(3.9 \times 10^{-3}/5 \times 10^{-7})/(14 - 4) = \ln 7800/10 = 0.896$।
\textbf{15.} $\ln 2/0.896 = 0.77$ समय-इकाई।
\textbf{16.} 0.906 से \qty{1}{\percent} कम: 200 जोड़ों का माध्य, एक सीमित समय और आलेख से लिए गए
पाठ्यांक।
\textbf{17.} जब दोनों रन दूर-दूर हो जाते हैं, तो वे आकर्षक के असंबद्ध बिंदुओं पर होते हैं, और
उनकी दूरी उसके आकार से अधिक नहीं हो सकती: वृद्धि संतृप्त हो जाती है।
\textbf{18.} $\ln 10/0.896 = 2.6$ समय-इकाई बाद।
\textbf{19.} $\log_2(3/0.1) = \log_2 30 = 4.91$ दोगुनीकरण।
\textbf{20.} $4.91 \times 2.4 = 11.8$~दिन।
\textbf{21.} एक दोगुना-काल: 2.4~दिन।
\textbf{22.} $4.91 \times 5 = 24.5$~दिन: पहले का अनुमान कहीं अधिक आशावादी था।
\textbf{23.} 2.4 दिनों के दोगुना-काल के साथ दस गुना छोटी त्रुटि $3.32 \times 2.4
= 8.0$~दिन देती, 4 से 5 नहीं: छोटी त्रुटियाँ 2.4 दिनों से कम में दोगुनी होती हैं, और उस
दोगुना-काल पर आधारित अनुमान क्षितिज को बढ़ा-चढ़ाकर बताता है।
\textbf{24.} यह पूर्वानुमान अधिक से अधिक लगभग \textbf{12 दिनों} तक उपयोगी हो सकता है, और व्यवहार में
इससे कम।
\end{solution}
